% U7 — KRoC 2027 단편 초안 (2쪽 목표; 공식 템플릿 공지 시 이식)
% 내용 축: 실패 주도 교육 시뮬레이터 소개 (U1 중심, U3와 축 분리)
\documentclass[10pt,a4paper,twocolumn]{article}
\usepackage{kotex}
\usepackage{amsmath,amssymb,bm}
\usepackage[margin=20mm]{geometry}
\usepackage{booktabs}
\usepackage{xurl}
\usepackage{hyperref}
\newcommand{\jj}{\mathrm{j}}

\title{\textbf{실패에서 배우는 임피던스 제어:\\실제 사고를 재현하는 공개 MuJoCo 교육 실험실}}
\author{정윤철 (Youncheol Jung)\\
\small 독립 연구자 | ORCID: 0009-0000-5282-881X}
\date{}

\begin{document}
\maketitle

\noindent\textbf{요약} ---
협동로봇의 보급으로 비전공 사용자도 임피던스 파라미터를 설정하는 시대가
오고 있으나, 강성과 목표 오프셋의 오설정은 고가 장비의 파손으로 직결된다.
본 논문은 저자가 실제 로봇에서 겪은 세 가지 사고 --- 무계획 파라미터 탐색,
저장 탄성 에너지 미고려로 인한 발사, 평형점 오해로 인한 목표 미도달 ---
를 안전하게 재현하는 공개 MuJoCo 교육 실험실을 소개한다.
1자유도 질량-스프링-댐퍼부터 7자유도 매니퓰레이터의 가상 벽 접촉까지
이어지는 4개 랩의 모든 설정에는 학습 가이드가 강제되며(테스트 스위트로
게이트), 원고의 수치 예제는 검증 스크립트로 기계 검산된다.

\section{서론}

임피던스 제어는 접촉 작업의 표준 언어이지만, 파라미터 설정의 직관은
문헌에서 얻기 어렵다.
유도가 생략된 논문과 상이한 표기 사이에서 초보자는 시행착오로 배우게
되는데, 실물에서 그 시행착오는 비용과 위험이 붙는 도박이다.
저자 역시 임피던스 설정 오류로 사물을 파손한 경험이 있다.
본 논문의 제안은 단순하다.
그 시행착오를 실물이 아니라 시뮬레이터에서, 실제 사고의 재현으로 겪게
하자는 것이다.

\section{실패 갤러리: 세 가지 사고}

공개 실험실\footnote{\url{https://github.com/ycpiglet/manipulator-control-tutorial}}의
실패 갤러리는 각 사고를 서사 → 원리 → 재현 설정 → 플롯 읽기 순서로
제공한다.

\textbf{사고 1: 무계획 파라미터 탐색.}
강성, 감쇠, 질량은 \(\zeta=d_d/(2\sqrt{m_dk_d})\)로 묶여 있어, 여러 값을
동시에 바꾸면 응답 변화의 원인을 분리할 수 없다.
실험실은 한 축씩만 바꾼 대비 설정 쌍(감쇠만, 강성만)을 제공해
one-factor-at-a-time 습관을 훈련시킨다.

\textbf{사고 2: 저장 탄성 에너지 발사.}
목표점에서 \(\delta\)만큼 떨어진 채 강성 \(k_d\)를 켜면
\(U=\frac{1}{2}k_d\delta^2\)이 장전되고, 무감쇠 근사에서 최고 속도는
\(v\approx\delta\sqrt{k_d/m}\)이다.
\(k_d=400~\mathrm{N/m}\), \(\delta=0.5~\mathrm{m}\), \(m=1~\mathrm{kg}\)이면
예측 \(10~\mathrm{m/s}\)이고, 시뮬레이터 측정값은 \(9.8~\mathrm{m/s}\)로
일치한다.
같은 이동거리를 \(k_d=30~\mathrm{N/m}\)로 명령하면 장전 에너지가 13배
작아진다.
명령 전 \(\frac{1}{2}k_d\delta^2\) 암산이 곧 장비 값이라는 것이 이 사례의
교훈이다.

\textbf{사고 3: 평형점 오해.}
임피던스 제어의 목표점 \(x_d\)는 위치 명령이 아니라 가상 스프링의
anchor이며, 접촉 시 로봇은 하중 평형점
\(x_{\mathrm{ss}}=x_d+f_{\mathrm{ext}}/k_d\)에서 멈춘다.
7자유도 Panda 가상 벽 실험에서 손끝은 벽면을 \(3.3~\mathrm{cm}\) 침투한
지점에서 멈추고, 측정 벽 힘 \(8.6~\mathrm{N}\)은 벽 강성
\(260~\mathrm{N/m}\times\)침투 깊이와 정확히 일치한다.
목표-실측 간극에 강성을 곱하면 접촉힘이 추정된다.

\section{실험실 구성과 품질 게이트}

실험실은 MuJoCo 기반 4개 랩(질량-스프링-댐퍼, PID, 2자유도 팔,
7자유도 Panda)으로 구성되며, 다음이 테스트로 강제된다.
(1) 모든 설정 파일은 학습 가이드(초점, 시도, 관찰, 다음 실행 제안)를
등록해야 하고, (2) 시뮬레이터 커버리지는 CI에서 80\% 하한이 강제되며
(현재 84\%), (3) 동반 원고의 수치 예제는 검증 스크립트의 수치
체크포인트로 기계 검산된다(오차 0).
코드는 Apache-2.0, 교육 콘텐츠는 CC BY 4.0로 공개되어 있다.

\section{결론}

비싼 시행착오를 무료의 재현 가능한 실험으로 바꾸는 것이 본 실험실의
기여이다.
발표에서는 세 사고의 라이브 재현과 파라미터 안전 설정 절차를 시연한다.
전체 이론(역사, 라플라스, 전기·기계 임피던스, 자코비안, 제어 방식
선택)을 유도 생략 없이 담은 확장판 원고가 같은 저장소에서 제공된다.

\end{document}
